\section{Related Work}
\textbf{Prompt Based Attacks.}
Prompt based attacks \citep{lee2024promptinfectionllmtollmprompt} exploit the LLMs by inserting a malicious or adversarial text into the user query. Prompt injection attacks \citep{liu2024promptinjectionattackllmintegrated} can be classified into two categories: (i) Direct Prompt Injection (DPI) and (ii) Indirect Prompt Injection (IPI) based on how the malicious instruction is injected. DPI involves embedding the malicious instruction directly into the user prompt to override or mislead the LLM into taking an action. \citet{ liu2024formalizingbenchmarkingpromptinjection} formalizes and benchmarks various instances of prompt injection involving diverse injected instructions. In contrast, IPI attacks \citep{zhan2024injecagentbenchmarkingindirectprompt, greshake2023youvesignedforcompromising, Yi_2025} rely on injecting the adversarial instruction into external sources such as tools, documents or web pages. The LLM then retrieves content from these external sources and incorporates the injected instruction, ultimately diverting the original user intent.  

\textbf{Agent Based Attacks.}
% \textcolor{red}{Complete this}
Agent-based attacks target the underlying system architecture rather than directly manipulating the LLMs themselves. For instance, \citet{wang2024badagentinsertingactivatingbackdoor} and \citet{Yang2024WatchOF} introduced a class of backdoor attacks where malicious triggers are embedded within the agent's environment, activating harmful behaviors when the agent accesses that environment. \citet{motwani2025secretcollusiongenerativeai} and \citet{wu2024shallteamupexploring} highlight the issue of colluding agents, a challenge particularly prevalent in multi-agent LLM frameworks where agents may collaborate toward a malicious goal. \cite{cemri2025multiagentllmsystemsfail} highlights failures arising from inter-agent misalignment and miscoordination.

\textbf{Safety evaluation of LLM Agents.}
% \textcolor{red}{Add details about AgentDojo, Inject Agent, R-judge, AgentMonitor, ASB, AgentHarm, AgentPoison, RedCode}
As LLM-based agents are increasingly deployed in real-world settings \citep{Xu2024TheAgentCompanyBL,Liu2023AgentBenchEL}, ensuring their safety and reliability has become a critical concern. Several benchmarks have been proposed to assess agent behavior under various adversarial and high-risk scenarios. AgentDojo \citep{debenedetti2024agentdojodynamicenvironmentevaluate} focuses on assessing prompt injection attacks and defenses, while InjectAgent \citep{zhan2024injecagentbenchmarkingindirectprompt} targets indirect prompt injection in contexts such as data security and financial harm. RedCode \citep{guo2024redcoderiskycodeexecution} benchmarks the ability of code agents to safely generate and execute potentially harmful code snippets. AgentHarm \citep{andriushchenko2025agentharmbenchmarkmeasuringharmfulness} evaluates how effectively agents refuse to comply with harmful or unethical queries. \cite{lee2024promptinfectionllmtollmprompt} study prompt propogation through self-replicating attacks while \cite{he-etal-2025-red} explore Agent-in-the-Middle attack. In contrast, R-Judge \citep{yuan2024rjudgebenchmarkingsafetyrisk} and AgentMonitor \citep{chan2024agentmonitorplugandplayframeworkpredictive} evaluate the safety awareness of LLMs by presenting them with manually curated records of risky agent trajectories, and assessing their ability to identify potential safety risks within those scenarios. 



